\section{Конечные автоматы}
	\tikzsetfigurename{RegLang_FA_}

\emph{Конечный автомат} (Finite State Machine, FSM)~--- вычислительная машина, которая имеет конечный набор состояний и может совершать переходы между ними, основываясь на прочитанных входных данных.
Важно отметить, что никакой дополнительной памяти классический конечный автомат не имеет и дополнительных действий (кроме чтения входной ленты) не производит%
\sidenote{Cуществуют автоматы с записью на отдельную (выходную) ленту (например, автоматы Мили~\cite{6771467}) и другие, основанные на конечных автоматах, вычислители, производящие дополнительные действия.
Например, расширенные конечные автоматы (Extended Finite State Machine), которые умеют работать с памятью (переменными)~\cite{Alagar2011, foster.ea:efsm:2020}.}.

Автоматом самого общего вида является \emph{недетерминированный конечный автомат}.
\begin{definition}[Недетерминированный конечный автомат]
    \label{def:NondeterminicticFiniteAutomata}
    \emph{Недетерминированный конечный автомат} (НКА, Nondeterministic Finite Automaton, NFA)~--- это пятёрка $M = \langle Q, Q_S, Q_F, \delta, \Sigma \rangle$, где
    \begin{itemize}
        \item $Q$~--- конечное множество состояний, $Q \neq \varnothing$;
        \item $Q_S \subseteq Q$~--- множество стартовых состояний, $Q_S \neq \varnothing$;
        \item $Q_F \subseteq Q$~--- множество финальных состояний\sidenote{Обратите внимание, множества стартовых и финальных состояний могут пересекаться.}, $Q_F \neq \varnothing$;
        \item $\delta \subseteq Q \times (\Sigma \cup \varepsilon) \times 2^Q$~--- функция переходов, а $\varepsilon \notin \Sigma$;
        \item $\Sigma$~--- конечный алфавит.
    \end{itemize}
\end{definition}

Для поиска путей в графе иногда требуется учитывать не только переходы по исходным меткам рёбер, но и переходы по рёбрам в обратном направлении.
Для этого удобно явно расширить алфавит обратными символами.
Пусть для каждого $a \in \Sigma$ задан символ $a^{-} \notin \Sigma$, причём $(a^{-})^{-} = a$.
Обозначим $\Sigma^{-} = \{a^{-} \mid a \in \Sigma\}$ и $\Sigma^{\leftrightarrow} = \Sigma \cup \Sigma^{-}$.

\mytodo{Дать более аккуратное определение 2-NFA из специализированной профильной литературы.}
\begin{definition}[Двусторонний недетерминированный конечный автомат]
    \label{def:TwoWayNondeterministicFiniteAutomata}
    \emph{Двусторонний недетерминированный конечный автомат} (Two-way Nondeterministic Finite Automaton, 2-NFA)~--- это недетерминированный конечный автомат $M = \langle Q, Q_S, Q_F, \delta, \Sigma^{\leftrightarrow} \rangle$, переходы которого помечены символами из расширенного алфавита $\Sigma^{\leftrightarrow}$.
    Переход по символу $a \in \Sigma$ соответствует чтению метки $a$ в прямом направлении, а переход по символу $a^{-} \in \Sigma^{-}$~--- чтению метки $a$ при движении по ребру графа в обратном направлении.
\end{definition}

Заметим, что функцию переходов можно представить разными способами: это может быть множество троек вида $(q_i, t, q_j)$, матрица, или же граф с метками на рёбрах
\[
G = \langle Q, \{(q_i,t,q_j \mid q_j \in \delta(q_i,t))\}, \Sigma \rangle.
\]
В дальнейшем мы чаще всего будем использовать представление автомата в виде графа.
Чтобы такое представление было полным, в графе отдельно обозначают стартовые и финальные состояния, как это показано на рисунке~\ref{fig:nfa_example}.

Так как нас интересуют конечные автоматы в контексте языков, то будем говорить, что на ленте автомата записано какое-то слово (или строка).
Иными словами, будем говорить, что автомат принимает на вход слово или строку.
При таком подходе автомат естественно рассматривать как \emph{распознаватель}\sidenote{Конечно, на автомат можно смотреть и как на генератор: любой путь от стартового состояния до финального задаёт слово из языка.}.

\begin{example}\label{exmpl:NFA}
    Пусть задан автомат $M=\langle Q, Q_S, Q_F, \delta, \Sigma\rangle$, где
    \begin{itemize}
        \item $Q = \{0,1,2\}$;
        \item $Q_S = \{0\}$;
        \item $Q_F = \{1\}$;
        \item $\Sigma = \{a,b\}$;
        \item $\delta$ определена следующим образом\sidenote{Часто указывают только те случаи, в которых $\delta$ возвращет непустое множество. Для неуказанных входных аргументов считается, что $\delta$ возвращает пустое множество.}:
        \begin{alignat*}{6}
            \delta(0,a) =& \{0,1\}      \qquad    & \delta(1,a) =& \varnothing \qquad & \delta(2,a) =& \varnothing \\
            \delta(0,b) =& \varnothing           & \delta(1,b) =& \{2\}             & \delta(2,b) =& \{1\} \\
            \delta(0,\varepsilon) =& \varnothing & \delta(1,\varepsilon) =& \{0\}   & \delta(2,\varepsilon) =& \varnothing.
        \end{alignat*}
    \end{itemize}
    Тогда его можно представить в виде графа как показано на рисунке~\ref{fig:nfa_example}.
    \begin{marginfigure}
        \begin{center}
            \input{part_02_Foundations/chapter_05_RegularLanguages/figures/02_FiniteAutomata/fig_nfa_example.tex}
        \end{center}
        \caption{Пример конечного автомата в котором состояние $0$~--- стартовое, а состояние $1$~--- финальное}
        \label{fig:nfa_example}
    \end{marginfigure}
\end{example}

Процесс вычислений, проделываемых конечным автоматом, удобно описывать в терминах переходов между \emph{конфигурациями}.

\begin{definition}[Конфигурация]
    Конфигурация $c$ конечного автомата $M = \langle Q, Q_S, Q_F, \delta, \Sigma \rangle$~--- это пара $(q, w)$, где $q\in Q$~--- это текущее состояние автомата, а $w \in \Sigma^*$~--- непросмотренная часть входной строки.
\end{definition}

\begin{definition}[Переход в НКА]
    Будем говорить, что автомат $M = \langle Q, Q_S, Q_F, \delta, \Sigma \rangle$ может перейти из конфигурации $c_1 = (q_1, w_1)$ в конфигурацию $c_2 = (q_2, w_2)$, если
    \[
    c_2 \in \{(q_2,w_2) \mid w_1 = aw_2, (q_1,a, q_2) \in \delta\} \cup \{(q_2,w_1) \mid (q_1, \varepsilon, q_2) \in \delta\}.
    \]
    Обозначать этот факт будем как $c_1 \to c_2$.
    Также будем считать, что на множестве конфигураций задано отношение перехода $(\to):(Q \times \Sigma^*)\times(Q \times \Sigma^*)$.

\end{definition}

\begin{definition}
    Транзитивное рефлексивное замыкание отношения перехода на конфигурациях будем обозначать следующим образом: \[ c_1 \to^* c_2. \]
    Альтернативно, в случае если $c_1 \to^* c_2$, будем говорить, что конфигурация $c_2$ \textit{достижима} из конфигурации $c_1$.
\end{definition}

Для удобства работы с недетерминированными автоматами расширим это отношение на множество конфигураций.

\begin{definition}
Будем говорить, что автомат может перейти из множества конфигураций $C_1$ в множество конфигураций $C_2$, если
\[
C_2 = \bigcup_{c_1 \in C_1} \{c_2 \mid c_1 \to c_2 \}.
\]

Обозначать этот факт будем как  $C_1 \Rightarrow C_2 $.
\end{definition}

\begin{definition}
Транзитивное рефлексивное замыкание отношения перехода на множествах конфигураций будем обозначать следующим образом: \[ C_1 \Rightarrow^* C_2. \]
\end{definition}

Для описания работы автомата $M = \langle Q, Q_S, Q_F, \delta, \Sigma \rangle$ нам понадобятся следующие выделенные типы конфигураций.
\begin{itemize}
    \item Стартовая конфигурация $c_s = (q_s,w)$, $q_s \in Q_S$, $w$~--- цепочка, которая подаётся на вход автомату.
    Для недетерминированного автомата естественно задать множество стартовых конфигураций $C_S = \bigcup_{q_s \in Q_S} (q_s,w)$.
    \item Финальная (принимающая) конфигурация $c_f = (q_f,\varepsilon)$, $q_f \in Q_F$.
\end{itemize}

Таким образом, работу автомата можно описать как последовательность переходов между множествами конфигураций.
Работа начинается с множества стартовых конфигураций и завершается в следующих двух случаях.
\begin{enumerate}
    \item Очередное множество конфигураций содержит финальную конфигурацию:
    \[C_S \Rightarrow^* C_i, c_f \in C_i.\] В этом случае говорят, что автомат \textit{принимает} входную строку.
    \item Очередное множество конфигураций пусто:
    \[C_0 = C_S \Rightarrow^* C_i \Rightarrow^* \varnothing, \text{ для любого } i: c_f \notin C_i.\]
    В этом случае говорят, что автомат \textit{не принимает} или \textit{отвергает} входную строку.
\end{enumerate}

\begin{definition}
    Язык задаваемый недетерминированным конечным автоматом $M = \langle Q, Q_S, Q_F, \delta, \Sigma \rangle$~--- это множество слов, принимаемых им:
     \[ \{w \mid (q_s,w) \to^* (q_f,\varepsilon), q_s \in Q_S, q_f \in Q_F \}.\]
\end{definition}

Так как конфигурация полностью описывает состояние процесса вычислений, то не надо обрабатывать одну и ту же конфигурацию несколько раз.
Это поможет при написании реального интерпретатора: будем отслеживать уже посещённые (обработанные) конфигурации\sidenote{Техника, аналогичная той, что применяется в обходах графов (обход в ширину, обход в глубину) для того, чтобы избежать повторного посещения вершин и, как следствие, зацикливания обхода. Более того, она типична для алгоритмов с рабочим множеством.}.

\begin{example}
    Рассмотрим пример работы конечного автомата, представленного на рисунке~\ref{fig:nfa_example}, на входной цепочке \texttt{abb}.
    В данном автомате стартовое состояние одно, потому множество стартовых конфигураций состоит из одной конфигурации:
    \[
    C_S = \{(0,\texttt{abb})\}.
    \]
    Из состояния \circled{0} существуют два перехода по символу \texttt{a}, который является первым в рассматриваемой цепочке.
    Потому будет получено следующее множество конфигураций:
    \[
        \{(0,\texttt{abb})\} \Rightarrow \{ (0,\texttt{bb}), (1,\texttt{bb})\}.
    \]
    В множестве появилась конфигурация, содержащая финальное состояние.
    Однако цепочка не прочитана до конца, потому продолжаем работу.
    Из нулевого состояния нет переходов по символу \texttt{b}, потому соответствующая конфигурация не породит новых.
    Из второй же конфигурации возможен переход по символу \texttt{b} и по $\varepsilon$:
    \[
        \{ (0,\texttt{bb}), (1,\texttt{bb})\} \Rightarrow \{ (0,\texttt{bb}), (2,\texttt{b})\}
    \]
    И снова, возможен только переход из состояния \circled{2}:
    \[
        \{ (0,\texttt{bb}), (2,\texttt{b})\} \Rightarrow \{(1,\varepsilon)\}.
    \]
    Теперь получена финальная конфигурация: состояние финальное и цепочка полностью обработана.
    Значит можно завершать работу и сообщать, что цепочка принимается автоматом.

    Полностью последовательность переходов выглядит следующим образом:
    \[
        \{(0,\texttt{abb})\} \Rightarrow \{ (0,\texttt{bb}), (1,\texttt{bb})\} \Rightarrow \{ (0,\texttt{bb}), (2,\texttt{b})\} \Rightarrow \{(1,\varepsilon)\}.
    \]
\end{example}

\begin{definition}[Детерминированный конечный автомат]
    \label{def:DeterminicticFiniteAutomata}
    \emph{Детерминированный конечный автомат} (ДКА, Deterministic Finite Automaton, DFA)~--- это пятёрка $M = \langle Q, q_S, Q_F, \delta, \Sigma \rangle$, где
    \begin{itemize}
        \item $Q$~--- конечное множество состояний, $Q \neq \varnothing$;
        \item $q_S \in Q$~--- стартовое состояние;
        \item $Q_F \subseteq Q$~--- множество финальных состояний\sidenote{Обратите внимание, стартовое состояние может входить в множество финальных.}\vspace{1.2\baselineskip}, $Q_F \neq \varnothing$;
        \item $\delta \subseteq Q \times \Sigma \times Q$~--- функция переходов\sidenote{При данном определении автомат будет ещё и \emph{полным}: для каждого состояния и каждого символа определён переход. Часто встречается вариант, когда функция переходов частично определена: для некоторых состояний не определены переходы для некоторых символов.};
        \item $\Sigma$~--- конечный алфавит.
    \end{itemize}
\end{definition}

Основных отличий от недетерминированного автомата два.
Во-первых, стартовое состояние ровно одно.
Во-вторых, функция переходов не допускает переходов по $\varepsilon$ и из любого состояния существует не более одного перехода по конкретному символу.

\begin{example}\label{exmpl:DFA}
    Пусть задан детерминированный автомат $M=\langle Q, q_S, Q_F, \delta, \Sigma\rangle$, где
    \begin{itemize}
        \item $Q = \{0,1,2\}$;
        \item $q_S = 0$;
        \item $Q_F = \{1\}$;
        \item $\Sigma = \{a,b\}$;
        \item $\delta$ определена следующим образом\sidenote{В нашем случае автомат не будет полным, так что $\delta$~--- частично определённая функция.}:
        \begin{alignat*}{6}
            \delta(0,a) =& 1      \qquad    & \delta(1,a) =& 1 \\
            \delta(1,b) =& 2      \qquad    & \delta(2,b) =& 1
        \end{alignat*}
    \end{itemize}
    Тогда его можно представить в виде графа, показанного на рисунке~\ref{fig:dfa_example}.
    \begin{marginfigure}
        \begin{center}
                        \begin{tikzpicture}
                \node[state,initial] (q_0) {$0$};
                \node[state,accepting] (q_1) [right = of q_0] {$1$};
                \node[state] (q_2) [above = of q_1] {$2$};
                \path[->]
                    (q_0) edge[bend left, above]   node {$a$} (q_1)
                    (q_1) edge[loop below, below]   node {$a$} (q_1)
                    (q_1) edge[bend left, left]   node {$b$} (q_2)
                    (q_2) edge[bend left, right]   node {$b$} (q_1);
            \end{tikzpicture}

        \end{center}
        \caption{Пример детерминированного конечного автомата в котором состояние $0$~--- стартовое, а состояние $1$~--- финальное, задающего тот же язык, что и автомат из примера~\ref{exmpl:NFA}}
        \label{fig:dfa_example}
    \end{marginfigure}

\end{example}

\begin{example}
    Рассмотрим пример работы конечного автомата, представленного на рисунке~\ref{fig:dfa_example}, на входной цепочке \texttt{abb}.
    Стартовая конфигурация:
    \[
    (0, \texttt{abb}).
    \]
    Из неё есть переход по символу \texttt{a}.
    Выполним его и получим следующую конфигурацию:
    \[
    (0, \texttt{abb}) \to (1, \texttt{bb}).
    \]
    Наблюдаем конфигурацию с финальным состоянием, но цепочка прочитана не до конца.
    Продолжаем работу:
    \[
    (1, \texttt{bb}) \to (2, \texttt{b}) \to (1, \varepsilon).
    \]
    Теперь конфигурация финальная: цепочка прочитана полностью и мы находимся в финально состоянии.
    Можно сообщить, что автомат принимает эту цепочку.
\end{example}
